% ======================================================================
\chapter[Ротатор. Вращательные и колебательные уровни молекул]%
        {Ротатор. Вращательные и колебательные уровни молекул}
\label{t:11}
% ======================================================================

\begin{qbox}
\textbf{Билет 11.} Уровни энергии плоского и пространственного ротатора. Вращательные и
колебательные уровни молекул.
\end{qbox}

\section{Содержание билета}

\textbf{Подтемы:}
\begin{enumerate}
 \item плоский ротатор: гамильтониан, уровни, волновые функции, вырождение;
 \item пространственный ротатор: уровни $E_J=\hbar^2J(J+1)/2I$, вырождение;
 \item двухатомная молекула: разделение электронного, колебательного и вращательного движений; оценки масштабов энергий;
 \item вращательный спектр (правило отбора, эквидистантные линии), колебательные уровни, ангармонизм, энергия диссоциации;
 \item колебательно-вращательный спектр (P- и R-ветви); проявления в теплоёмкости газов.
\end{enumerate}

\textbf{Необходимые понятия:} оператор момента импульса и его собственные значения
(\tref{t:09}); гармонический осциллятор (\tref{t:08}); приведённая масса (\tref{t:03}).

\section{Определения и физические принципы}

\begin{definition}[ротатор]
\emph{Ротатор} — система, способная только вращаться вокруг неподвижной точки (центра масс) с
фиксированным моментом инерции $I$. \emph{Плоский} ротатор вращается в фиксированной плоскости
(одна угловая координата $\varphi$); \emph{пространственный} — вращается произвольно (углы $\theta,\varphi$).
Модель «жёсткой гантели»: две массы $m_1$, $m_2$ на расстоянии $r_0$, $I=\mu r_0^2$, $\mu=m_1m_2/(m_1+m_2)$.
\end{definition}

\begin{definition}[вращательная постоянная]
\begin{equation}
  B=\frac{\hbar^2}{2I}\qquad\Bigl(\text{в спектроскопии также } \tilde B=\frac{B}{hc}=\frac{\hbar}{4\pi cI}\ \text{в см}^{-1}\Bigr).
\end{equation}
\end{definition}

\emph{Физический смысл.} Классически энергия вращения $E=L^2/2I$ непрерывна. В квантовой
механике $L^2$ квантуется, поэтому квантуется и энергия вращения. \emph{Область применимости:}
модель жёсткого ротатора справедлива, когда центробежное растяжение связи мало
($BJ(J+1)\ll\hbar\omega_{\text{кол}}$) и вращение не возбуждает электронные переходы.

\section{Основная теория}

\subsection{Плоский ротатор}

Кинетическая энергия вращения в плоскости $E=L_z^2/2I$, поэтому гамильтониан
\begin{equation}\label{eq:11:Hplane}
  \op H=\frac{\op L_z^2}{2I}=-\frac{\hbar^2}{2I}\diff{^2}{\varphi^2}.
\end{equation}
Собственные функции $\op L_z$ (\tref{t:09}) — одновременно собственные функции $\op H$
(вывод~\ref{der:11:plane}):
\begin{equation}\label{eq:11:Eplane}
  \Phi_m=\frac{\ee^{\ii m\varphi}}{\sqrt{2\pi}},\qquad E_m=\frac{\hbar^2m^2}{2I},\qquad m=0,\pm1,\pm2,\dots
\end{equation}
Уровни с $m\neq0$ \emph{двукратно вырождены} ($\pm m$ — вращение в противоположных направлениях с одинаковой
энергией); основное состояние $m=0$ невырождено и имеет $E=0$ — нулевых «вращательных колебаний» нет,
поскольку координата $\varphi$ циклическая и ограничения на неё нет (в отличие от ямы или осциллятора).

\subsection{Пространственный ротатор}

\begin{equation}\label{eq:11:Hsph}
  \op H=\frac{\op L^2}{2I}\quad\Longrightarrow\quad
  E_J=\frac{\hbar^2}{2I}J(J+1)=BJ(J+1),\qquad J=0,1,2,\dots,
\end{equation}
волновые функции — сферические функции $Y_{JM}(\theta,\varphi)$, кратность вырождения $g_J=2J+1$
($M=-J,\dots,J$) — энергия не зависит от ориентации оси вращения. Для молекул вращательное квантовое
число обозначают $J$. Расстояние между соседними уровнями растёт линейно:
\begin{equation}\label{eq:11:dE}
  E_{J+1}-E_J=2B(J+1).
\end{equation}

\subsection{Двухатомная молекула: иерархия движений}

Молекула — система из ядер (масса $M$) и электронов (масса $\mel\ll M$). Лёгкие электроны движутся
быстро, ядра — медленно; поэтому в нулевом приближении (приближение Борна--Оппенгеймера) электронная
задача решается при неподвижных ядрах, а её энергия $U(r)$ как функция межъядерного расстояния
служит потенциалом для движения ядер (рис.~\ref{fig:11:morse}). Полная энергия
\begin{equation}\label{eq:11:Etot}
  E=E_{\text{эл}}+E_{\text{кол}}+E_{\text{вр}},\qquad
  E_{\text{эл}}:E_{\text{кол}}:E_{\text{вр}}\sim1:\sqrt{\frac{\mel}{M}}:\frac{\mel}{M}
\end{equation}
(вывод~\ref{der:11:scales}). Численно: $E_{\text{эл}}\sim$ единицы эВ (видимый и УФ диапазоны),
$E_{\text{кол}}\sim0{,}1$--$0{,}5$~эВ (инфракрасный), $E_{\text{вр}}\sim10^{-4}$--$10^{-2}$~эВ (далёкий ИК и микроволновый).

\subsection{Колебательные уровни}

Вблизи минимума $r_0$ потенциальной кривой $U(r)\approx U(r_0)+\frac12k(r-r_0)^2$. Относительное движение ядер
вдоль оси молекулы — осциллятор с приведённой массой $\mu$:
\begin{equation}\label{eq:11:vib}
  E_v=\hbar\omega\Bigl(v+\frac12\Bigr),\qquad \omega=\sqrt{\frac{k}{\mu}},\qquad v=0,1,2,\dots
\end{equation}
Реальный потенциал ангармоничен (например, потенциал Морзе $U=D_e(1-\ee^{-\beta(r-r_0)})^2$): с ростом $v$
уровни сгущаются,
\begin{equation}\label{eq:11:anh}
  E_v=\hbar\omega\Bigl(v+\frac12\Bigr)-\hbar\omega x_e\Bigl(v+\frac12\Bigr)^2,\qquad x_e\ll1,
\end{equation}
и сходятся к границе диссоциации. \emph{Энергия диссоциации} из основного состояния
$D_0=D_e-\hbar\omega/2$ меньше глубины ямы $D_e$ на энергию нулевых колебаний. Правило отбора для гармонических
колебаний $\Delta v=\pm1$; ангармонизм допускает слабые обертоны $\Delta v=\pm2,\pm3,\dots$ Колебательный
дипольный спектр есть только у полярных молекул (HCl, CO); у гомоядерных (H$_2$, N$_2$) дипольный момент не
меняется при колебаниях, и ИК-поглощения нет (колебания наблюдаются в комбинационном рассеянии).

\begin{figure}[htbp]
\centering
\begin{tikzpicture}
\begin{axis}[width=0.75\linewidth, height=6.0cm, axis lines=left, xlabel={$r$}, ylabel={$U(r)$},
   xmin=0.4, xmax=4.6, ymin=-0.05, ymax=1.25, xtick={1.2}, xticklabels={$r_0$}, ytick={0,1},
   yticklabels={$0$,$D_e$}, samples=300, clip=false]
  \addplot[very thick, domain=0.735:4.5] {(1-exp(-1.6*(x-1.2)))^2};
  \addplot[gray, dashed, domain=0.4:4.6] {1};
  \foreach \v in {0,1,2,3,4,5,6} {
    \pgfmathsetmacro{\Ev}{0.16*(\v+0.5)-0.0128*(\v+0.5)^2}
    \pgfmathsetmacro{\xl}{1.2-ln(1+sqrt(\Ev))/1.6}
    \pgfmathsetmacro{\xr}{1.2-ln(1-sqrt(\Ev))/1.6}
    \edef\tmp{\noexpand\draw[accent, thick] (axis cs:\xl,\Ev) -- (axis cs:\xr,\Ev);}\tmp
  }
  % вращательные подуровни для v=0 (схематично)
  \foreach \J in {1,2,3} { \pgfmathsetmacro{\EJ}{0.0768+0.0035*\J*(\J+1)}
     \edef\tmp{\noexpand\draw[keycol] (axis cs:1.05,\EJ) -- (axis cs:1.40,\EJ);}\tmp }
  \node[font=\small, accent, anchor=west] at (axis cs:2.35,0.36) {колебательные уровни $v$};
  \node[font=\small, anchor=west] at (axis cs:2.6,0.9) {граница диссоциации};
  \draw[<->] (axis cs:0.55,0) -- (axis cs:0.55,1) node[midway, left, font=\small] {$D_e$};
  \draw[<->, dercol] (axis cs:4.3,0.0768) -- (axis cs:4.3,1) node[midway, right, font=\small] {$D_0$};
  \draw[dercol, dotted] (axis cs:1.4,0.0768) -- (axis cs:4.3,0.0768);
\end{axis}
\end{tikzpicture}
\caption{Потенциальная кривая двухатомной молекулы (потенциал Морзе, схема): колебательные уровни
сгущаются к границе диссоциации; $D_0=D_e-\hbar\omega/2$. Около уровня $v=0$ условно показаны вращательные
подуровни (в реальных молекулах они значительно гуще).}
\label{fig:11:morse}
\end{figure}

\subsection{Вращательный спектр}

Переходы между вращательными уровнями одного колебательного состояния возможны для молекул с
постоянным дипольным моментом; правило отбора $\Delta J=\pm1$. Частоты поглощения $J\to J+1$:
\begin{equation}\label{eq:11:rotspec}
  \hbar\omega_J=E_{J+1}-E_J=2B(J+1),\qquad J=0,1,2,\dots
\end{equation}
— эквидистантная последовательность линий с интервалом $2B$ (рис.~\ref{fig:11:spec}). По интервалу
определяют $I$ и межъядерное расстояние $r_0=\sqrt{I/\mu}$. Интенсивности линий определяются
заселённостями $N_J\propto(2J+1)\ee^{-BJ(J+1)/k_{\mathrm B}T}$ — максимум при $J\approx\sqrt{k_{\mathrm B}T/2B}$.

\begin{figure}[htbp]
\centering
\begin{tikzpicture}[xscale=1.25]
  \draw[->] (0,0) -- (9.8,0) node[right] {$\omega$};
  \foreach \J/\h in {0/0.9,1/1.6,2/2.0,3/2.05,4/1.8,5/1.4,6/0.95,7/0.55} {
     \draw[very thick, accent] ({1.1*(\J+1)},0) -- ({1.1*(\J+1)},\h);
  }
  \draw[<->] (1.1,-0.35) -- (2.2,-0.35) node[midway, below, font=\small] {$2B/\hbar$};
  \node[font=\small] at (1.1,2.4) {$0\!\to\!1$};
  \node[font=\small] at (3.3,2.4) {$2\!\to\!3$};
\end{tikzpicture}
\caption{Вращательный спектр поглощения полярной двухатомной молекулы: эквидистантные линии
$\hbar\omega=2B(J+1)$; огибающая интенсивностей определяется больцмановской заселённостью
уровней (схема).}
\label{fig:11:spec}
\end{figure}

\subsection{Колебательно-вращательный спектр}

При переходе $v=0\to1$ одновременно меняется $J$ ($\Delta J=\pm1$ для двухатомных молекул в электронном состоянии $^1\Sigma$,
к которому относится основное состояние большинства двухатомных молекул). При одинаковой $B$ в обоих колебательных состояниях:
\begin{align}
  \text{R-ветвь }(J\to J+1):\quad &\hbar\omega=\hbar\omega_0+2B(J+1),\quad J=0,1,\dots\\
  \text{P-ветвь }(J\to J-1):\quad &\hbar\omega=\hbar\omega_0-2BJ,\quad J=1,2,\dots
\end{align}
Спектр — две серии линий с интервалом $2B$ по обе стороны от $\omega_0$, а на месте $\omega_0$ — провал
(переход $\Delta J=0$ запрещён). Из-за того что во втором колебательном состоянии $r$ в среднем больше,
$B_1<B_0$ — линии R-ветви сгущаются, P-ветви — разрежаются. Изотопический эффект: у H$^{35}$Cl и H$^{37}$Cl
различаются $\mu$, поэтому линии раздвоены ($\omega\propto\mu^{-1/2}$).

\paragraph{Характерные значения.}
\begin{center}
\small
\begin{tabular}{@{}lccc@{}}
\toprule
Молекула & $r_0$, Å & $B$, мэВ & $\hbar\omega$, эВ \\
\midrule
H$_2$ & 0{,}74 & 7{,}5 & 0{,}55 \\
HCl & 1{,}27 & 1{,}3 & 0{,}37 \\
CO & 1{,}13 & 0{,}24 & 0{,}27 \\
\bottomrule
\end{tabular}
\end{center}

\subsection{Теплоёмкость двухатомного газа}

Степени свободы «включаются», когда $k_{\mathrm B}T$ становится сравнимой с расстоянием между уровнями. Для H$_2$
характерные температуры $\Theta_{\text{вр}}=B/k_{\mathrm B}\approx88$~К, $\Theta_{\text{кол}}=\hbar\omega/k_{\mathrm B}\approx6300$~К.
Поэтому при $T\ll\Theta_{\text{вр}}$ $C_V=\frac32R$ (только поступательное движение), при
$\Theta_{\text{вр}}\ll T\ll\Theta_{\text{кол}}$ $C_V=\frac52R$, и лишь при $T\gtrsim\Theta_{\text{кол}}$ $C_V\to\frac72R$.
Классическая теория (равнораспределение) не могла объяснить «вымерзание» степеней свободы.

\section{Математические выводы}

\begin{derivation}[label=der:11:plane]{уровни плоского ротатора}
Стационарное уравнение: $-\dfrac{\hbar^2}{2I}\dfrac{\dd^2\Phi}{\dd\varphi^2}=E\Phi$, т.\,е. $\Phi''+\dfrac{2IE}{\hbar^2}\Phi=0$.
Общее решение $\Phi=A\ee^{\ii\kappa\varphi}+B\ee^{-\ii\kappa\varphi}$, $\kappa=\sqrt{2IE}/\hbar$. Условие однозначности
$\Phi(\varphi+2\pi)=\Phi(\varphi)$ даёт $\ee^{2\pi\ii\kappa}=1$, т.\,е. $\kappa=|m|$, $m$ — целое. Отсюда
\[
  E_m=\frac{\hbar^2m^2}{2I},\qquad \Phi_m=\frac{1}{\sqrt{2\pi}}\ee^{\ii m\varphi}.
\]
Тот же результат сразу следует из $\op H=\op L_z^2/2I$ и $L_z=m\hbar$: собственная функция $\op L_z$ — собственная
функция $\op L_z^2$ с собственным значением $m^2\hbar^2$. Каждому $E_m$ ($m\neq0$) отвечают две независимые функции
$\ee^{\pm\ii|m|\varphi}$ — двукратное вырождение.
\end{derivation}

\begin{derivation}[label=der:11:rigid]{ротатор как частица на сфере}
Пусть частица массы $\mu$ движется по сфере радиуса $r_0$ (жёсткая гантель после отделения центра масс).
Гамильтониан $\op H=-\frac{\hbar^2}{2\mu}\nabla^2$ при фиксированном $r=r_0$ содержит лишь угловую часть лапласиана:
\[
  \op H=-\frac{\hbar^2}{2\mu r_0^2}\Delta_{\theta,\varphi}=\frac{\op L^2}{2\mu r_0^2}=\frac{\op L^2}{2I}.
\]
Собственные значения $\op L^2$ равны $\hbar^2J(J+1)$ (\tref{t:09}), откуда $E_J=\hbar^2J(J+1)/2I$, $g_J=2J+1$.
Классический предел: при $J\gg1$ $E\approx\hbar^2J^2/2I=L^2/2I$ с $L\approx\hbar J$.
\end{derivation}

\begin{derivation}[label=der:11:scales]{оценка соотношения электронных, колебательных и вращательных энергий}
\emph{Электронная энергия.} Электрон локализован в области размера $a\sim\aB$:
$E_{\text{эл}}\sim\hbar^2/(\mel a^2)\sim$ несколько эВ.

\emph{Колебательная энергия.} Сдвиг ядер на расстояние $\sim a$ меняет электронную энергию на величину
порядка $E_{\text{эл}}$, поэтому жёсткость связи $k\sim E_{\text{эл}}/a^2$. Частота колебаний ядер массы $M$:
\[
  \hbar\omega=\hbar\sqrt{\frac kM}\sim\hbar\sqrt{\frac{E_{\text{эл}}}{Ma^2}}
  =\hbar\sqrt{\frac{\hbar^2}{\mel a^2}\cdot\frac{1}{Ma^2}}=\frac{\hbar^2}{\mel a^2}\sqrt{\frac{\mel}{M}}
  =E_{\text{эл}}\sqrt{\frac{\mel}{M}} .
\]
\emph{Вращательная энергия.} $I\sim Ma^2$, $E_{\text{вр}}\sim\dfrac{\hbar^2}{Ma^2}=\dfrac{\hbar^2}{\mel a^2}\dfrac{\mel}{M}=E_{\text{эл}}\dfrac{\mel}{M}$.

При $\mel/M\sim10^{-4}$ (лёгкие молекулы): $E_{\text{эл}}:E_{\text{кол}}:E_{\text{вр}}\sim1:10^{-2}:10^{-4}$.
Эта иерархия оправдывает раздельное рассмотрение трёх типов движения. Попутно: амплитуда нулевых
колебаний $\delta r\sim\sqrt{\hbar/M\omega}\sim a(\mel/M)^{1/4}\sim0{,}1a$ — ядра действительно локализованы.
\end{derivation}

\begin{derivation}[label=der:11:H2]{вращательная постоянная H$_2$}
$\mu=m_p/2$, $r_0=0{,}741$~Å. $I=\mu r_0^2$; в энергетических единицах
\[
  B=\frac{\hbar^2}{2\mu r_0^2}=\frac{(\hbar c)^2}{2\mu c^2r_0^2}=\frac{(1973\ \text{эВ}\cdot\text{Å})^2}{2\cdot469\cdot10^6\ \text{эВ}\cdot(0{,}741\ \text{Å})^2}
  =\frac{3{,}89\cdot10^6}{5{,}15\cdot10^8}\ \text{эВ}\approx7{,}6\cdot10^{-3}\ \text{эВ}.
\]
Интервал между линиями вращательного спектра $2B\approx15$~мэВ соответствует $\lambda=hc/2B\approx0{,}08$~мм
(далёкий ИК). Температура $\Theta_{\text{вр}}=B/k_{\mathrm B}=7{,}6\cdot10^{-3}/8{,}62\cdot10^{-5}\approx88$~К.
\end{derivation}

\begin{derivation}[label=der:11:PR]{P- и R-ветви}
Энергия состояния $E(v,J)=\hbar\omega_0(v+\frac12)+BJ(J+1)$. Переход $(0,J)\to(1,J')$:
\[
  \Delta E=\hbar\omega_0+B\bigl[J'(J'+1)-J(J+1)\bigr].
\]
При $J'=J+1$: $J'(J'+1)-J(J+1)=(J+1)(J+2)-J(J+1)=2(J+1)$ — R-ветвь.
При $J'=J-1$: $(J-1)J-J(J+1)=-2J$ — P-ветвь. При $J'=J$ было бы $\Delta E=\hbar\omega_0$, но такой
переход запрещён правилом отбора $\Delta J=\pm1$ (фотон уносит момент $\hbar$), поэтому в центре спектра — провал.
\end{derivation}

\section{Итоговые формулы}

\keyformula{E_m=\frac{\hbar^2m^2}{2I},\ \ m=0,\pm1,\dots;\qquad E_J=\frac{\hbar^2}{2I}J(J+1),\ \ g_J=2J+1}
 {$I=\mu r_0^2$ — момент инерции, $m$ — проекция момента (плоский ротатор), $J$ — вращательное квантовое число.}
 {жёсткий ротатор (плоский / пространственный).}
 {вращательные уровни молекул.}

\keyformula{\hbar\omega=2B(J+1),\qquad B=\frac{\hbar^2}{2I},\qquad \Delta J=\pm1}
 {$B$ — вращательная постоянная.}
 {полярные молекулы, переходы внутри одного колебательного состояния.}
 {по интервалу $2B$ между линиями — момент инерции и межъядерное расстояние.}

\keyformula{E_v=\hbar\omega\Bigl(v+\frac12\Bigr)-\hbar\omega x_e\Bigl(v+\frac12\Bigr)^2,\qquad \omega=\sqrt{k/\mu},\qquad D_0=D_e-\frac{\hbar\omega}{2}}
 {$v$ — колебательное число, $x_e$ — постоянная ангармонизма, $D_e$ — глубина ямы, $D_0$ — энергия диссоциации.}
 {малый ангармонизм.}
 {колебательный спектр, энергия диссоциации, изотопические эффекты.}

\keyformula{E_{\text{эл}}:E_{\text{кол}}:E_{\text{вр}}\sim1:\sqrt{\mel/M}:\mel/M}
 {$\mel$ — масса электрона, $M$ — масса ядер.}
 {оценка по порядку величины.}
 {обоснование разделения движений; спектральные диапазоны (УФ-видимый, ИК, микроволновый).}

\section{Физическая интерпретация}

\textbf{Структура молекулярных спектров.} Каждый электронный уровень содержит систему колебательных,
а каждый колебательный — систему вращательных подуровней. Поэтому электронные переходы в молекулах
дают не линии, а \emph{полосы} (полосатые спектры), а ИК-переходы — колебательно-вращательные полосы.

\textbf{Зависимость от параметров.} $B\propto1/(\mu r_0^2)$ — тяжёлые молекулы имеют гуще расположенные
уровни; $\omega\propto\sqrt{k/\mu}$ — изотопическое замещение (H$\to$D) уменьшает $\omega$ в $\approx\sqrt2$ раз
и $B$ — в 2 раза (для H$_2$/D$_2$).

\textbf{Связь с термодинамикой.} Квантование объясняет вымерзание вращательных и колебательных
степеней свободы и температурную зависимость теплоёмкости газов — одно из ранних подтверждений
квантовой теории.

\section{Что необходимо знать на экзамене}

\begin{examcheck}
\checkhead{Определения}
\begin{itemize}
 \item плоский и пространственный ротатор, момент инерции, вращательная постоянная; приближение Борна--Оппенгеймера.
\end{itemize}
\checkhead{Формулы}
\begin{itemize}
 \item $E_m=\hbar^2m^2/2I$ (вырождение 2); $E_J=BJ(J+1)$ (вырождение $2J+1$);
 \item $\hbar\omega=2B(J+1)$, $\Delta J=\pm1$; $E_v=\hbar\omega(v+\frac12)$ (+ангармонизм); $D_0=D_e-\hbar\omega/2$;
 \item P- и R-ветви; $E_{\text{эл}}:E_{\text{кол}}:E_{\text{вр}}\sim1:\sqrt{m/M}:m/M$.
\end{itemize}
\checkhead{Выводы}
\begin{itemize}
 \item уровни плоского ротатора из условия однозначности; пространственный ротатор через $\op L^2$;
 \item оценки масштабов энергий; численная оценка $B$ для H$_2$; ветви колебательно-вращательного спектра.
\end{itemize}
\checkhead{Интерпретации}
\begin{itemize}
 \item почему у ротатора нет нулевой энергии, а у осциллятора есть;
 \item спектральные диапазоны; полосатые спектры; вымерзание степеней свободы и теплоёмкость.
\end{itemize}
\end{examcheck}
